\section{问题四的建模与求解}

问题四在前三问基础上进一步引入逐日波动电价，沿用微网购电与储能基础模型，分别在波动电价条件下重新求解问题二和问题三。


\subsection{电价预测与动态更新}

附件4给出了2025年各日不同时段的实际电价，但在制定日前购电计划时，
目标日未来电价尚不可知。因此，利用目标日前已经获得的历史电价构造预测值，
并比较不同历史信息选取方法。

首先，考虑不同星期类型的电价变化规律。将星期一至星期日分别定义为七种星期类型，
选取目标日前与其星期类型相同的最近两个历史日，以对应时段电价均值作为目标日预测电价。
设$H_d$为第$d$日之前与其星期类型相同的最近至多两个历史日集合，
则第$d$日第$t$个时段的预测电价为
\begin{equation}
\widehat{p}_{d,t}
=
\frac{1}{|H_d|}
\sum_{h\in H_d}p_{h,t}.
\label{eq:p4_price_wd2}
\end{equation}

当同星期历史数据不足时，采用前一日对应时段的实际电价；
若尚无可用历史电价，则以附件1中的分时电价作为初始值。

为比较不同预测方式，进一步设置两种对照方法。
其一，直接采用前一日对应时段的实际电价：
\begin{equation}
\widehat{p}_{d,t}=p_{d-1,t}.
\label{eq:p4_price_yesterday}
\end{equation}
其二，直接采用附件1中的固定分时电价。
后续分别比较三种方法的电价预测误差和购电费用。

在波动电价下重新求解问题三时，除更新光伏发电功率预报外，
还利用调整时刻之前已经获得的实际电价修正剩余时段预测电价。
设$\tau$为当前调整时刻，$t_0$为第一个尚未执行时段，
$\widehat{p}^{(0)}_{d,t}$为当天$0{:}00$得到的初始预测电价，
根据已执行时段实际电价与初始预测电价的平均水平计算修正系数
\begin{equation}
r_d^{(\tau)}
=
\frac{
\displaystyle \frac{1}{t_0-1}\sum_{t=1}^{t_0-1}p_{d,t}
}{
\displaystyle \frac{1}{t_0-1}\sum_{t=1}^{t_0-1}\widehat{p}^{(0)}_{d,t}
}.
\label{eq:p4_price_ratio}
\end{equation}

据此更新尚未执行时段的预测电价：
\begin{equation}
\widehat{p}^{(\tau)}_{d,t}
=
r_d^{(\tau)}\widehat{p}^{(0)}_{d,t},
\qquad t\geq t_0.
\label{eq:p4_price_update}
\end{equation}

已经执行时段采用实际电价，尚未执行时段采用更新后的预测电价，
保证日内调整仅使用当前时刻已经获得的信息。


\subsection{波动电价下问题二的模型}

在问题二模型基础上，将固定分时电价替换为预测电价，
负载与光伏发电预测方法、储能运行约束及跨日状态衔接方式保持不变。
每天$0{:}00$根据预测负载、预测光伏发电量和预测电价制定全天计划，
目标函数为
\begin{equation}
\min C_d^{\mathrm{plan}}
=
\sum_{t=1}^{144}
\left(
\widehat{p}_{d,t}x_{d,t}
+
5\widehat{p}_{d,t}\widehat{e}_{d,t}
\right),
\label{eq:p4_2_plan_cost}
\end{equation}
其中，$x_{d,t}$表示第$d$天第$t$个时段的计划购电量，
$\widehat{e}_{d,t}$表示预测条件下的供电缺口。

计划确定后，各时段计划购电量和储能充放电量在当天保持不变。
根据目标日实际负载和实际光伏发电量计算实际供电缺口：
\begin{equation}
e_{d,t}
=
\max
\left\{
0,\,
L_{d,t}^{\mathrm{act}}
+u_{d,t}
-x_{d,t}
-R_{d,t}^{\mathrm{act}}
-v_{d,t}
\right\}.
\label{eq:p4_2_emergency}
\end{equation}

紧急购电按照交易时刻实际电价的5倍计价，因此当天总购电费用为
\begin{equation}
C_d^{4\text{-}2}
=
\sum_{t=1}^{144}p_{d,t}x_{d,t}
+
5\sum_{t=1}^{144}p_{d,t}e_{d,t}.
\label{eq:p4_2_actual_cost}
\end{equation}

由此得到波动电价下问题二的日前购电及储能调度策略。


\subsection{波动电价下问题三的模型}

在问题三日内调整模型基础上进一步引入预测电价及动态更新。
每天$0{:}00$根据负载预测、附件3中的未来24 h光伏发电功率预报
以及初始预测电价制定计划购电策略：
\begin{equation}
\min C_d^{(0)}
=
\sum_{t=1}^{144}
\left(
\widehat{p}^{(0)}_{d,t}x^{\mathrm{plan}}_{d,t}
+
5\widehat{p}^{(0)}_{d,t}\widehat{e}^{(0)}_{d,t}
\right).
\label{eq:p4_3_plan}
\end{equation}

其中，$x^{\mathrm{plan}}_{d,t}$为$0{:}00$制定的计划购电量，
$\widehat{p}^{(0)}_{d,t}$为初始预测电价，
$\widehat{e}^{(0)}_{d,t}$为预测供电缺口。

在$6{:}00$、$12{:}00$和$18{:}00$获得新的光伏发电功率预报后，
仅对尚未执行的时段进行调整。
设$\mathcal{T}_{\tau}$为调整时刻$\tau$后的未执行时段集合，
$d^-_{d,t}$和$d^+_{d,t}$分别表示相对于原计划减少和增加的购电量，则
\begin{equation}
x^{(\tau)}_{d,t}
=
x^{\mathrm{plan}}_{d,t}
-d^-_{d,t}
+d^+_{d,t},
\qquad t\in\mathcal{T}_{\tau}.
\label{eq:p4_3_adjust}
\end{equation}

按照题目规定，减少计划购电量需支付交易时刻电价50\%的违约费用，
增加部分按照交易时刻电价的1.5倍计价。
结合紧急购电费用，剩余时段优化目标为
\begin{equation}
\min C_d^{(\tau)}
=
\sum_{t\in\mathcal{T}_{\tau}}
\left[
\widehat{p}^{(\tau)}_{d,t}x^{(\tau)}_{d,t}
+
0.5\widehat{p}^{(\tau)}_{d,t}
\left(d^-_{d,t}+d^+_{d,t}\right)
+
5\widehat{p}^{(\tau)}_{d,t}\widehat{e}^{(\tau)}_{d,t}
\right].
\label{eq:p4_3_adjust_cost}
\end{equation}

考虑5倍电价紧急购电的高成本，正式策略仅允许在原计划基础上增加购电量：
\begin{equation}
d^-_{d,t}=0,\qquad d^+_{d,t}\geq0,
\qquad t\in\mathcal{T}_{\tau}.
\label{eq:p4_3_no_reduction}
\end{equation}

同时保持$0{:}00$确定的储能充放电方案不变，
利用新增光伏预报和更新后的电价信息调整尚未执行时段的购电量。
每次重新求解后，分别计算沿用上一轮方案和候选调整方案的预计费用，
仅当候选方案费用更低时接受调整。
所有调整量及相关费用均以$0{:}00$原计划为基准计算。

当天运行结束后，采用附件4中的实际电价计算最终费用。
设$x^{\mathrm{act}}_{d,t}$为最终执行购电量，
$e^{\mathrm{act}}_{d,t}$为实际紧急购电量，则
\begin{equation}
C_d^{4\text{-}3}
=
\sum_{t=1}^{144}
\left[
p_{d,t}x^{\mathrm{act}}_{d,t}
+
0.5p_{d,t}
\left|
x^{\mathrm{act}}_{d,t}
-x^{\mathrm{plan}}_{d,t}
\right|
\right]
+
5\sum_{t=1}^{144}
p_{d,t}e^{\mathrm{act}}_{d,t}.
\label{eq:p4_3_actual_cost}
\end{equation}


\subsubsection{滚动调整策略设置}

为比较电价更新和储能调整方式，设置四种滚动策略，
其余模型参数和约束保持一致。

\begin{table}[H]
\centering
\caption{波动电价下不同滚动购电策略设置}
\label{tab:p4_3_strategy}
\resizebox{\textwidth}{!}{
\begin{tabular}{ccccc}
\toprule
策略 & 新光伏预报 & 电价信息处理 & 储能重新优化 & 主要考察内容 \\
\midrule

A
& 不利用
& 不进行日内更新
& 否
& 无日内调整的基准策略 \\

B0
& 利用
& 采用前一日电价，不进行日内更新
& 否
& 基础滚动调整 \\

bstar
& 利用
& 同星期历史日均值法（WD2）并动态更新
& 是
& 储能同步调整 \\

bfreeze
& 利用
& 同星期历史日均值法（WD2）并动态更新
& 否
& 冻结储能计划滚动调整 \\

\bottomrule
\end{tabular}}
\end{table}

策略A不利用日内新增信息；
B0采用前一日电价并调整尚未执行时段购电量；
bstar在更新购电量的同时重新优化储能；
bfreeze保持$0{:}00$储能充放电方案不变，仅调整购电量。
通过比较紧急购电量和总购电费用确定最终策略。


\subsection{统计期策略比较}

对2025年2月1日至12月31日共334天进行逐日求解。
对于问题二，比较不同电价预测方法；
对于问题三，进一步比较不同日内调整方式。


\subsubsection{波动电价下问题二的比较}

保持问题二的负载预测、光伏发电预测及储能运行方式不变，
仅改变计划阶段采用的电价预测方法。
三种方法的统计期结果如表~\ref{tab:p4_2_price_compare}所示。

\begin{table}[H]
\centering
\caption{不同电价预测方法下问题二的统计期计算结果}
\label{tab:p4_2_price_compare}
\resizebox{\textwidth}{!}{
\begin{tabular}{lrrrrrr}
\toprule
方法 &
电价预测MAE（元/kWh） &
计划购电量（kWh） &
计划购电费用（元） &
紧急购电量（kWh） &
紧急购电费用（元） &
总购电费用（元） \\
\midrule

附件1固定电价法（Att1）
& 0.096323
& 22,556,016.294
& 14,580,647.832
& 318,742.269
& 1,283,945.585
& \textbf{15,864,593.417} \\

前一日电价法（Yest）
& 0.084326
& 22,595,609.093
& 14,610,113.777
& 320,839.031
& 1,290,713.507
& 15,900,827.284 \\

同星期历史日均值法（WD2）
& \textbf{0.045153}
& \textbf{22,528,946.052}
& \textbf{14,566,044.707}
& 325,100.369
& 1,305,493.684
& 15,871,538.391 \\
\bottomrule
\end{tabular}}
\end{table}

同星期历史日均值法（WD2）的电价预测MAE最低，为$0.045153$元/kWh，
较前一日电价法和附件1固定电价法分别降低约$46.45\%$和$53.12\%$，
其计划购电量和计划购电费用也均为三种方法中的最低值。

但WD2的统计期总购电费用为$15,871,538.391$元，
比附件1固定电价法高$6,944.974$元。
电价预测MAE只衡量预测价格与实际价格之间的偏差，
最终购电费用还受到负载与光伏发电预测误差、储能调度及紧急购电发生时段影响。
因此，以电价预测误差最低作为选择依据，后续采用WD2作为主要动态电价预测方法，
并保留附件1固定电价法作为对照。这一选择侧重电价预测精度，并不意味着WD2在三种方法中实现了最低的实际总购电费用。


\subsubsection{波动电价下问题三的比较}

在波动电价条件下重新求解问题三，
比较不进行日内调整、采用前一日电价滚动调整、
储能同步调整和冻结储能计划滚动调整四种策略，
结果见表~\ref{tab:p4_3_strategy_compare}。

\begin{table}[H]
\centering
\caption{不同滚动调整策略下问题三的统计期计算结果}
\label{tab:p4_3_strategy_compare}
\resizebox{\textwidth}{!}{
\begin{tabular}{llrrrrrrr}
\toprule
策略 &
价格方法 &
计划购电量 &
最终执行购电量 &
购电及调整费用 &
紧急购电量 &
紧急购电费用 &
总购电费用 &
接受调整次数 \\
&
&
（kWh） &
（kWh） &
（元） &
（kWh） &
（元） &
（元） &
\\
\midrule

A：不调整
& WD2
& 22,732,368.289
& 22,732,368.289
& 14,683,855.828
& 390,699.520
& 1,617,838.822
& 16,301,694.650
& 0 \\

B0：前一日电价滚动调整
& Yest
& 22,803,070.535
& 23,477,152.769
& 15,572,980.794
& \textbf{151,679.212}
& \textbf{638,123.936}
& 16,211,104.730
& 805 \\

bstar：储能同步调整
& WD2
& 22,732,368.289
& 23,241,290.868
& 15,272,017.972
& 283,029.339
& 1,204,478.958
& 16,476,496.930
& 814 \\

bfreeze：冻结储能计划滚动调整
& WD2
& 22,732,368.289
& 23,405,782.648
& 15,528,793.583
& 153,358.381
& 643,835.663
& \textbf{16,172,629.246}
& 805 \\
\bottomrule
\end{tabular}}
\end{table}

bfreeze的统计期总购电费用最低，为$16,172,629.246$元。
相较不调整策略A，紧急购电量由$390,699.520$ kWh降至
$153,358.381$ kWh，下降约$60.75\%$；
总购电费用由$16,301,694.650$元降至$16,172,629.246$元，
减少$129,065.404$元，降幅约$0.79\%$。

B0的紧急购电量为$151,679.212$ kWh，略低于bfreeze，
但总购电费用高出$38,475.484$元；
bstar同步重新优化储能后，总费用达到$16,476,496.930$元。
综合紧急购电风险与总购电费用，最终选取bfreeze作为波动电价下
问题三的滚动调整策略。


\subsection{月度与指定日期比较}

为考察bfreeze在不同时期的运行表现，将其与不调整策略A进行比较。
两种策略均采用WD2进行初始电价预测，月度结果见表~\ref{tab:p4_monthly_compare}。

\begin{table}[H]
\centering
\caption{不同月份滚动调整策略的计算结果}
\label{tab:p4_monthly_compare}
\resizebox{\textwidth}{!}{
\begin{tabular}{crrrrr}
\toprule
月份 &
A总购电费用（元） &
bfreeze总购电费用（元） &
费用差额（元） &
A紧急购电量（kWh） &
bfreeze紧急购电量（kWh） \\
\midrule
2025-02 & 1,443,700.859 & 1,396,963.281 &  46,737.578 & 48,617.178 & 27,932.062 \\
2025-03 & 1,320,304.649 & 1,264,500.244 &  55,804.405 & 51,229.194 & 24,921.974 \\
2025-04 & 1,157,316.814 & 1,071,342.329 &  85,974.484 & 69,389.861 & 26,441.428 \\
2025-05 & 1,118,022.695 & 1,105,981.340 &  12,041.355 & 40,888.984 & 10,482.229 \\
2025-06 & 1,639,260.724 & 1,672,057.517 & -32,796.793 & 26,289.413 &  8,608.481 \\
2025-07 & 1,934,717.286 & 1,906,909.177 &  27,808.109 & 58,560.922 & 23,866.106 \\
2025-08 & 1,605,946.447 & 1,600,027.316 &   5,919.132 & 34,527.412 & 14,360.746 \\
2025-09 & 1,289,031.926 & 1,300,641.353 & -11,609.427 & 17,492.893 &  2,719.820 \\
2025-10 & 1,295,027.579 & 1,293,403.942 &   1,623.637 & 19,832.367 &  6,738.333 \\
2025-11 & 1,447,387.096 & 1,491,741.831 & -44,354.735 &  7,162.512 &  1,447.334 \\
2025-12 & 2,050,978.574 & 2,069,060.917 & -18,082.343 & 16,708.784 &  5,839.866 \\
\bottomrule
\end{tabular}}
\end{table}

滚动调整在11个月中均降低了紧急购电量。
其中2月、3月、4月、5月、7月、8月和10月的总购电费用下降，
6月、9月、11月和12月有所增加。
当新增购电及调整费用超过紧急购电费用的减少时，
总购电费用仍可能增加。



进一步比较题目指定的2025年3月20日、6月21日、9月23日和12月21日，
结果见表~\ref{tab:p4_typical_days}。

\begin{table}[H]
\centering
\caption{指定日期不同滚动调整策略的计算结果}
\label{tab:p4_typical_days}
\resizebox{\textwidth}{!}{
\begin{tabular}{crrrrrr}
\toprule
日期 &
A总购电费用（元） &
bfreeze总购电费用（元） &
费用差额（元） &
A紧急购电量（kWh） &
bfreeze紧急购电量（kWh） &
接受调整次数 \\
\midrule
2025-03-20 & 56,112.558 & 52,458.706 &  3,653.851 & 3,129.442 & 1,958.551 & 2 \\
2025-06-21 & 20,384.342 & 20,937.865 &   -553.523 & 1,027.296 &   488.916 & 3 \\
2025-09-23 & 53,436.261 & 53,674.350 &   -238.088 &    12.514 &    12.514 & 2 \\
2025-12-21 & 76,855.778 & 77,782.925 &   -927.147 &   409.862 &   203.647 & 2 \\
\bottomrule
\end{tabular}}
\end{table}

3月20日采用bfreeze后，紧急购电量由$3,129.442$ kWh降至
$1,958.551$ kWh，总购电费用减少$3,653.851$元。
6月21日和12月21日的紧急购电量分别下降约$52.41\%$和$50.31\%$，
但总费用分别增加$553.523$元和$927.147$元；
9月23日紧急购电量保持$12.514$ kWh，总费用增加$238.088$元。

\subsection{问题四购电策略与结果}

问题四在逐日波动电价条件下分别重新求解问题二和问题三，并据此确定相应购电策略。

对于问题二，保持原有负载与光伏发电预测、储能运行及跨日状态衔接方式不变，以预测电价替代固定分时电价。以电价预测精度为选择依据，采用预测MAE最低的\textbf{同星期历史日均值法（WD2）}，其MAE为$0.045153$元/kWh。该策略在统计期内计划购电量为$22,528,946.052$ kWh，计划购电费用为$14,566,044.707$元，紧急购电量为$325,100.369$ kWh，\textbf{实际总购电费用为$15,871,538.391$元}。作为对照，附件1固定电价法的总购电费用为$15,864,593.417$元，略低于WD2，表明\textbf{电价预测MAE较低并不必然使实际购电费用同步降低}。


表~\ref{tab:p4_2_selected_days}进一步列出波动电价下问题二采用WD2时四个指定日期的结果，购电量与费用直接取自\texttt{result4-2.xlsx}。其中计划购电费用不包含紧急购电费用，紧急购电量按当日全部非零记录汇总。

\begin{table}[H]
\centering
\caption{波动电价下问题二四个指定日期的购电结果（WD2）}
\label{tab:p4_2_selected_days}
\begin{tabular}{crrr}
\toprule
日期 & 计划购电量（kWh） & 计划购电费用（元） & 紧急购电量（kWh） \\
\midrule
2025-03-20 & 71494.161 & 44772.588 & 302.639 \\
2025-06-21 & 36076.246 & 17858.487 & 362.215 \\
2025-09-23 & 73017.301 & 47883.122 & 1079.032 \\
2025-12-21 & 98134.741 & 73856.078 & 881.389 \\
\bottomrule
\end{tabular}
\end{table}

四个日期中，6月21日的计划购电量与计划购电费用最低，12月21日最高；但紧急购电量最高的是9月23日，说明计划购电规模与预测偏差造成的实际供电缺口并不呈简单对应关系。上述结果分别反映日前计划与日终缺口结算，不能仅用计划费用判断最终总费用。

对于问题三，每天$0{:}00$根据负载预测、光伏发电功率预报和WD2预测电价制定初始购电及储能方案；在$6{:}00$、$12{:}00$和$18{:}00$获得新光伏预报后，利用当日已获得的实际电价更新剩余时段预测电价，并调整尚未执行的购电量。策略比较后采用\textbf{冻结$0{:}00$储能充放电方案的bfreeze策略}，统计期紧急购电量为$153,358.381$ kWh，\textbf{总购电费用为$16,172,629.246$元}；相较不调整策略，紧急购电量\textbf{下降$60.75\%$}，总购电费用\textbf{减少$129,065.404$元}。因此，在波动电价条件下，\textbf{利用日内新增光伏预报和已获得的实际电价调整剩余购电计划，能够降低统计期紧急购电量和总购电费用}。

波动电价下问题二、问题三的完整逐时段购电策略、储能充放电量及紧急购电结果见附件。